En esta práctica se observó qué le pasa a un grupo de rayos de luz paralelos cuando atraviesan una pieza de acrílico de forma cilíndrica, es decir, cuando cruzan dos superficies refractantes curvas. Con una fuente de luz y una máscara de varias rendijas se obtuvieron siete rayos paralelos que se hicieron incidir sobre la cara curva del cilindro. Fue una práctica de observación, sin mediciones. Al salir del acrílico los rayos dejaron de ser paralelos y se juntaron en una zona muy cercana a la pieza, para luego volver a separarse. Se notó además que los rayos de las orillas se cruzaban antes que los del centro, de modo que no había un único punto de enfoque, y que en los bordes de los rayos aparecían colores. Estas observaciones coinciden con la ecuación de una superficie refractante esférica, que para un cilindro de acrílico predice un foco a una distancia de aproximadamente la mitad del radio detrás de la pieza, y con los fenómenos conocidos como aberración esférica y dispersión.
